\section{重心}\label{節:面積重心と線重心}

前節の重み付き平均を,密度が一様な薄い板の重心に使ってみよう.板を小さな部分に分け,\,$i$ 番目の部分の面積を $\Delta A_i$,位置をほぼ $(x_i,y_i)$ とする.密度が一様なら各部分の質量は面積に比例し,共通の密度は分子と分母で約分できる.よって重心の座標は,面積を重みとした平均
\[
\overline x\approx\frac{\sum_i x_i\Delta A_i}{\sum_i\Delta A_i},\qquad
\overline y\approx\frac{\sum_i y_i\Delta A_i}{\sum_i\Delta A_i}
\]
で表される.小さな部分として縦の細い帯をとり,この和を定積分に帰着させよう.

$a<b$ とし,\,$f,g$ は $[a,b]$ で連続で $g(x)\le f(x)$ を満たすとする.板が
\[
a\le x\le b,\qquad g(x)\le y\le f(x)
\]
で表され,その面積 $A$ は正であるとする.位置 $x$ の幅 $\Delta x$ の帯は,面積がほぼ $\{f(x)-g(x)\}\Delta x$ の長方形で,その中心は $\left(x,\frac{f(x)+g(x)}{2}\right)$ である.\,$y$ 座標の重み付きの和を作ると,帯1つ分は
\[
\frac{f(x)+g(x)}{2}\cdot\{f(x)-g(x)\}\Delta x
 =\frac12\{f(x)^2-g(x)^2\}\Delta x
\]
となる.分割を細かくすると
\begin{equation}\label{式:面積の重心}
\begin{aligned}
A&=\int_a^b(f-g)\dd x,\\
\overline x&=\frac1A\int_a^b x(f-g)\dd x,\\
\overline y&=\frac1{2A}\int_a^b(f^2-g^2)\dd x
\end{aligned}
\end{equation}
を得る.これを\textbf{面積重心}という.\,$\overline y$ を計算するとき,上端の高さ $f$ だけを帯の面積に掛けてはいけない.帯の重心は上端ではなく,上下の中央にある.

曲線に沿って曲げた一様な針金では,質量が比例するのは面積ではなく長さである.曲線 $C$ の長さを $L>0$ とし,曲線の小さな部分の長さを $\dd s$ で表すと,重心は
\begin{equation}\label{式:線の重心}
\overline x=\frac1L\int_C x\dd s,\qquad
\overline y=\frac1L\int_C y\dd s
\end{equation}
となる.これを\textbf{線重心}という.\,$\int_C y\dd s$ は新しい種類の積分ではなく,曲線を媒介変数で表せば,式\eqref{式:媒介曲線の弧長}と同じく
\[
\int_C y\dd s
 =\int_\alpha^\beta y(t)\sqrt{x'(t)^2+y'(t)^2}\dd t
\]
という普通の定積分で計算する.針金の重心を求めるときは,曲線を1度だけたどる区間を選ぶ.

\begin{bookpoint}
\textbf{POINT}\par
板では縦の帯の面積 $(f-g)\dd x$ を重みにして,帯の中心 $\left(x,\frac{f+g}{2}\right)$ の座標を平均する.針金では長さ $\dd s$ を重みにする.
\end{bookpoint}

\paragraphbf{半円板と半円弧}
半径 $a$ の上半円板と,その円弧だけでできた針金の重心を比べよう.どちらも $y$ 軸について対称なので $\overline x=0$ である.

半円板では $A=\frac{\pi a^2}{2}$,\,$g(x)=0$,\,$f(x)=\sqrt{a^2-x^2}$ である.式\eqref{式:面積の重心}の $\overline y$ の積分は
\begin{align*}
\frac12\int_{-a}^a(a^2-x^2)\dd x
 &=\frac12\left[a^2x-\frac{x^3}{3}\right]_{-a}^a\\
 &=\frac{2a^3}{3}
\end{align*}
となる.これを $A$ で割って,半円板の重心の高さは
\begin{equation}\label{式:半円板の重心}
\overline y_{\mathrm{面}}=\frac{2a^3}{3}\cdot\frac{2}{\pi a^2}=\frac{4a}{3\pi}
\end{equation}
である.

円弧は $x=a\cos t$,\,$y=a\sin t$,\,$0\le t\le\pi$ と表せる.\,$\sqrt{x'(t)^2+y'(t)^2}=a$ より $\dd s=a\dd t$,\,$L=\pi a$ であるから,
\begin{equation}\label{式:半円弧の重心}
\overline y_{\mathrm{線}}
 =\frac{1}{\pi a}\int_0^\pi(a\sin t)a\dd t
 =\frac{2a}{\pi}
\end{equation}
となる.\,$\frac{2a}{\pi}\approx0.64a$,\,$\frac{4a}{3\pi}\approx0.42a$ なので,円弧だけのほうが重心は高い.円板には直径の近くにも面積があるが,円弧だけの針金にはその部分がないからである.
